\scanpage{28}
\mspage{9}
with $E_J\cap T$ conical in $T$, and the conical structure
\deleted{induced} being the product structure by $I^J$ of
a conical structure on $(T_J,Y_J)$, having
$\dot T_J$ as boundary-part of the cone on $Y_J$\(\mathinner{\cdotp\cdotp-}\).

\deleted{\unsure{6}}

\textcircled{6.} Relations to stratifications of type $I$
($I$ finite ordered set)
Such a stratification (of a space $X$) is
the datum of an increasing family of
closed subspaces $X_i$. We say that
the stratification is \uline{precylindrical} if

\deleted{$\forall i$}
a) $\forall j'\leftrightarrow j$, $j$ successor of $j'$,
$X_{j'}\hookrightarrow X_j$ is
\deleted{\unsure{a subset}} \inserted{bounding}

\inserted{b) \deleted{\unsure{For}}\inserted{For} $j$ fixed,
$j'$ varying among the predecessors, the
$X_{j'}$ are transverse \unsure{bounding} subsets of $X_j$}

\margin{i.e. equipped with the transverse cylindrical neighborhoods}

\deleted{b)}\inserted{c)}
$\forall i\leftrightarrow j$, the structure of the set \inserted{$I_{i,j}$}
of $k$ such that $i\le k\le j$ is that of the set
\inserted{of subsets of a set} \deleted{simplex} --- which means that
if $S\inserted{(=S_{ij})}$ is the set of predecessors of $j$ in
$I_{i,j}$ \inserted{i.e. of the minimal elements of $I_{i,j}\setminus\{j\}$}
and if for \unsure{$\forall$} $k\in I_{i,j}$, we denote
\deleted{\unsure{then}} by $s_k$ \deleted{the set of}\inserted{the} \unsure{subsets|subset} $s_k\subset S$ formed
of $j\in S_{\unsure{k}}$ such that
$j\ge s_k$ (thus $s_k=s_{kj}$)
the map $k\longrightarrow s_k$ (which is decreasing)
is an anti-isomorphism of $I_{ij}$ with $\mathcal P(S)$
ordered by inclusion. \deleted{$\unsure{I}$}

\margin{\uline{NB} Taking c) into account, it suffices to require in
b) for $j$ maximal \unsure{in $I$.}}

\margin{\deleted{\unsure{in one and the same cone}}
\inserted{\deleted{\unsure{subsets}}}}

\margin{We further want \deleted{\unsure{that}}
that the infima in $I_{i,j}$ correspond\unsure{} to
the intersection of the $X_k$\(\mathinner{\cdotp-\, -}\) cf. (0) below.}

Combining b) and c), we see that the
\margin{27}
